% Generated draft: requires mathematical and source review.
\documentclass[11pt]{article}
\usepackage[margin=25mm]{geometry}
\usepackage{fontspec}
\setmainfont{Latin Modern Roman}
\usepackage{amsmath,amssymb,mathtools,unicode-math}
\setmathfont{Latin Modern Math}
\usepackage{enumitem,needspace,xurl,hyperref}
\setlength{\parindent}{0pt}
\setlength{\parskip}{5pt}
\setlength{\emergencystretch}{3em}
\newcommand{\N}{\mathbb{N}}
\newcommand{\Z}{\mathbb{Z}}
\newcommand{\Q}{\mathbb{Q}}
\newcommand{\R}{\mathbb{R}}
\newcommand{\C}{\mathbb{C}}
\newcommand{\F}{\mathbb{F}}
\newcommand{\A}{\mathbb{A}}
\newcommand{\PP}{\mathbb{P}}
\newcommand{\E}{\mathbb{E}}
\newcommand{\eps}{\varepsilon}
\newcommand{\dd}{\,\mathrm{d}}
\begin{document}
\section*{UnsolvedMath ID: 6700067}
\textbf{Source problem code:} catalogue entry
\subsection{Definitions and mathematical statement}
Let $B^n(R)$ be the Euclidean radius-$R$ ball and let $S$ be the intended metric target. Define
\[
\operatorname{Lip}_1(B^n(R),S)=\{f:B^n(R)\to S:d_S(f(x),f(y))\leq|x-y|\text{ for all }x,y\}.
\]
Assume $H_i(S;\mathbb R)=0$ for $i=1,\ldots,n$, with real coefficients. The source asks whether
\[
D_c(R)=\operatorname{diam}_c\bigl(\operatorname{Lip}_1(B^n(R),S)\bigr)
\]
stays bounded as $R\to\infty$ for a sufficiently large fixed $c$. A basic uniform interpretation is that suitable finite $c,C,R_0$, depending on the fixed target and dimension but not on $R$, satisfy $D_c(R)\leq C$ for all $R\geq R_0$. The parameter $c$ must not grow with $R$.

The subscripted diameter is not defined in the supplied record. One possible mapping-space length convention would use the uniform distance $d_\infty(f,g)=\sup_x d_S(f(x),g(x))$, allow homotopies through $c$-Lipschitz maps, and minimize their $d_\infty$-length between endpoint maps. This is explanatory only: the source may use a different path norm, and it does not say whether disconnected pairs have infinite distance or diameters are taken componentwise. The geometry and connectedness of $S$, ball-boundary conventions, and the exact quantifier on sufficiently large $c$ are also absent. Thus the boundedness target is retained without inventing the missing diameter definition.
\subsection{Short English statement}
If the target has vanishing real homology through degree $n$, do the source's controlled mapping-space diameters of $1$-Lipschitz maps from larger and larger $n$-balls stay uniformly bounded with fixed control $c$? The controlled distance is not defined in the extract.
\subsection{Sources}
\begin{itemize}
\item Misha Gromov, \emph{101 Questions, Problems and Conjectures around Scalar Curvature} (2017), printed tag [?71], occurrence 70, PDF page 72, subpart [a], \url{https://www.ihes.fr/~gromov/wp-content/uploads/2018/08/101-problemsOct1-2017.pdf}; accessed 2026-07-29, PDF-text-normalized extraction. The provenance describes the public author PDF as incomplete.
\item The supplied summary associates this with Gromov's mapping-space path-length ideas and refers to ``Length of paths'' without a complete publication identifier or the required distance definition.
\item Local registry \texttt{data/unsolvedmath/problems.json}, permanent UnsolvedMath ID $6700067$, source code catalogue entry, \emph{Scalar Curvature Question [?71]}. Dataset attribution: Ulam AI and the UnsolvedMath Contributors.
\end{itemize}
\end{document}
